% Рисунок с доски: схема вязко-упругого материала Кельвина — Фойгта:
% пружина (1) параллельно демпферу (2).
\begin{figure}[htbp]
\centering
\begin{tikzpicture}[>={Latex[length=4pt]}]
  \pic[rotate=90] at (0,0) {spring};
  \node[anchor=east, font=\footnotesize] at (-0.45,0) {$1$};
  \node[anchor=east, font=\footnotesize] at (-0.45,0.5) {$\mu$};

  \pic[rotate=90] at (1.8,0) {damper};
  \node[anchor=west, font=\footnotesize] at (2.3,0) {$2$};
  \node[anchor=west, font=\footnotesize] at (2.3,0.5) {$\eta$};

  \draw[thick] (0,0.9) -- (0,1.3) -- (1.8,1.3) -- (1.8,0.9);
  \draw[thick] (0,-0.9) -- (0,-1.3) -- (1.8,-1.3) -- (1.8,-0.9);
  \draw[thick] (0.9,1.3) -- (0.9,1.8);
  \draw[thick] (0.9,-1.3) -- (0.9,-1.8);

  \node[anchor=west, align=left] at (3.4,0.6)
    {$\Sdev = \Sdev_1 + \Sdev_2$};
  \node[anchor=west, align=left] at (3.4,-0.4)
    {$\dev = \dev_1 = \dev_2$};
\end{tikzpicture}
\caption{Материал Кельвина–Фойгта: пружина и демпфер параллельно. Отсюда
$\Sdev = 2\mu\,\dev + 2\eta\,\dt{\dev}$.}
\label{fig:kelvin-voigt}
\end{figure}
